\documentclass[11pt]{article}

% ---------------------------------------------------------------------------
% Pump / gain-map continuation methods: current scheme and alternatives.
% Standalone; compiles with a stock TeX Live (pdflatex).
% ---------------------------------------------------------------------------

\usepackage[a4paper,margin=1in]{geometry}
\usepackage{amsmath,amssymb}
\usepackage{booktabs}
\usepackage{array}
\usepackage{xcolor}
\usepackage{listings}
\usepackage{enumitem}
\usepackage[hidelinks]{hyperref}

\newcommand{\code}[1]{\texttt{\small #1}}
\newcommand{\Res}{\mathbf{R}}
\newcommand{\Xv}{\mathbf{X}}
\newcommand{\Jac}{\mathbf{J}}

\lstset{
  basicstyle=\ttfamily\footnotesize,
  columns=fullflexible,
  breaklines=true,
  frame=single,
  framesep=4pt,
  backgroundcolor=\color{black!3},
  keywordstyle=\color{blue!60!black},
  commentstyle=\color{green!40!black},
}

\title{Traversing the Pump/Gain Map:\\
Continuation and Warm-Start Strategies in \code{twpa\_solver}\\
\large Current Scheme and a Survey of Alternatives}
\author{TWPA solver --- \code{twpa\_jax}}
\date{\today}

\begin{document}
\maketitle

\begin{abstract}
The gain map produced by \code{scripts/run\_gain\_map.py} requires a converged
harmonic-balance (HB) pump solution at every cell of a two-dimensional
(pump-power $\times$ pump-frequency) grid. Because each cell is a stiff
nonlinear root-find, the map is not solved cell-by-cell from scratch; instead
the solver \emph{continues} solutions from one operating point to the next.
This report documents exactly which continuation and warm-start strategies the
code uses today, locates each in the source, and then surveys more sophisticated
alternatives from the numerical-continuation literature (pseudo-arclength,
tangent predictors, deflated continuation, pseudo-transient continuation, and
2-D map-traversal schemes), with references. It is descriptive: it maps the
current design and catalogues options, and does not prescribe a change.
\end{abstract}

\tableofcontents

% ===========================================================================
\section{Problem setup}
% ===========================================================================

\subsection{The harmonic-balance residual}
Each map cell fixes a pump frequency $f_p$ (angular $\omega_p = 2\pi f_p$) and a
pump drive amplitude $I_p$. The pump state is the vector of complex harmonic
coefficients $\Xv$ of the node fluxes on the mode set
$\{k\,\omega_p\}$. The steady state solves a nonlinear algebraic residual
\begin{equation}
  \Res(\Xv;\ \lambda) \;=\; 0,
  \qquad
  \Res(\Xv;\lambda) \;=\; \mathcal{L}(\omega_p)\,\Xv \;+\; \mathcal{N}(\Xv) \;-\; \lambda\, \mathbf{S},
  \label{eq:residual}
\end{equation}
where $\mathcal{L}(\omega_p)$ is the frequency-domain linear block, $\mathcal{N}$
the Josephson (cosine) nonlinearity evaluated through the alias-free time
transform, $\mathbf{S}$ the pump source phasor, and
$\lambda \in [0,1]$ a \emph{scalar homotopy / continuation parameter} that
scales the drive. At $\lambda = 1$ the physical problem is recovered; at
$\lambda \to 0$ the problem is linear and trivially solvable. The source enters
\emph{linearly} in $\lambda$ (Eq.~\ref{eq:residual}), which is the property the
whole continuation machinery exploits (\code{src/twpa\_solver/pump/solver.py}).

\subsection{Two nested traversal problems}
There are two distinct ``traversal'' questions, and the code answers them with
different mechanisms:
\begin{enumerate}[leftmargin=*]
  \item \textbf{Intra-cell (parameter homotopy).} How do we get \emph{one} cell
    from the easy linear regime up to full drive $\lambda=1$? --- answered by
    \emph{continuation in $\lambda$} (\S\ref{sec:intracell}).
  \item \textbf{Inter-cell (map traversal).} How do we use a \emph{neighbouring}
    converged cell to cheapen the next cell? --- answered by \emph{warm-starting}
    plus an extrapolation predictor along the power axis (\S\ref{sec:intercell}).
\end{enumerate}

% ===========================================================================
\section{What we do now}
% ===========================================================================

\subsection{Intra-cell: natural-parameter continuation in the pump amplitude}
\label{sec:intracell}

The continuation parameter is the pump drive amplitude itself (via $\lambda$),
i.e. \emph{natural-parameter continuation}. Two modes exist.

\paragraph{Fixed-step continuation (the ``cold'' reference).}
\code{solve\_continuation} (\code{solver.py:402}) walks a fixed ladder
$\lambda \in \{1/N, 2/N, \dots, 1\}$ with $N=20$ steps
(\code{--continuation-steps}, default 20). At each step it runs one
Newton--Krylov solve seeded from the previous converged step. If any step fails
it returns immediately (no re-convergence is attempted). This is the trusted,
deterministic path used to validate everything else and is selected by
\code{pump\_flags\_cold} (\code{run\_gain\_map.py:252}).

\paragraph{Adaptive-step continuation (the ``seed'' path).}
\code{solve\_adaptive\_continuation} (\code{solver.py:463}) replaces the fixed
ladder with a step controller: start at \code{--adaptive-initial-step}, on a
converged step grow the step by $\times 1.5$, on a failed step shrink by
$\times 0.5$; if the step falls below \code{--adaptive-min-step} it falls back to
a fixed 20-step run. This is the first-point path for each frequency column in
the warm pass, seeded from the linear phasor guess (below).

\paragraph{The Newton--Krylov corrector.}
Every continuation step calls \code{solve\_one} (\code{solver.py:101}): a damped
Newton iteration whose linear systems are solved matrix-free by GMRES against
the true Jacobian $\Jac = \partial\Res/\partial\Xv$, preconditioned by one of
several factorizations (\code{--inproc-preconditioner}). Cost-control knobs live
here: a stall detector (\code{stall\_ratio}, \code{stall\_patience}) and a
wall-clock \code{solve\_deadline\_s} abort stiff \emph{over-fold} solves, and
modified-Newton preconditioner reuse (\code{precond\_reuse}) amortizes the LU.

\paragraph{Predictors inside continuation.}
At each $\lambda$ step the initial guess is, by default, simply the previous
converged state (a \emph{trivial / zeroth-order predictor}). Optionally a
\emph{secant predictor} extrapolates linearly in $\lambda$ from the last two
converged states (\code{solver.py:422} and \code{:501}):
\begin{equation}
  \Xv_{\text{guess}} \;=\; \Xv_{\text{prev}}
    \;+\; \beta\,\bigl(\Xv_{\text{prev}} - \Xv_{\text{prevprev}}\bigr),
  \qquad
  \beta = \frac{\lambda - \lambda_{\text{prev}}}
               {\lambda_{\text{prev}} - \lambda_{\text{prevprev}}} .
  \label{eq:secant}
\end{equation}

\paragraph{The linear phasor seed (homotopy start).}
\code{build\_linear\_phasor\_seed} (\code{seeds.py:46}) solves the
\emph{first-harmonic linear block only}, $\mathcal{L}_1(\omega_p)\,\Xv_1 =
\mathbf{S}_1$, by GMRES, giving a physically-motivated non-zero start for the
adaptive path. This is the $\lambda \to 0$ analytic solution used as the
continuation origin (as opposed to the all-zeros cold start).

\subsection{Inter-cell: warm-starting across the map}
\label{sec:intercell}

The map is \emph{not} solved cell-independently. The production path is the
in-process warm pass \code{run\_warm\_pass\_inprocess}
(\code{run\_gain\_map.py:849}), which traverses the grid as follows:

\begin{enumerate}[leftmargin=*]
  \item \textbf{Column-major, increasing power.} Cells are grouped by frequency
    column (\code{by\_col}); each column is sorted by ascending pump power and
    traversed low$\to$high (\code{run\_gain\_map.py:871--880}).
  \item \textbf{First cell of a column: seed.} Solved from the
    \code{linear\_phasor} seed with adaptive continuation (mode \code{"seed"}).
  \item \textbf{Subsequent cells: warm start.} Each higher-power cell is solved
    by a \emph{single full-scale Newton solve at $\lambda=1$},
    \code{solve\_direct} (\code{solver.py:383}), initialised from the previous
    converged cell's HB state --- \emph{no} intra-cell continuation ladder. This
    is the main speed-up: one Newton solve instead of $\sim 20$.
  \item \textbf{Secant predictor along the power axis.} Between cells the guess
    is extrapolated from the last two converged solutions along the
    pump-\emph{current} axis, \code{secant\_guess} (\code{run\_gain\_map.py:785},
    default \code{--inproc-fold-predictor secant}). Structurally identical to
    Eq.~\ref{eq:secant} but with the pump current as the abscissa. A bad
    extrapolation is caught and retried from the plain warm start
    (\code{:901}).
  \item \textbf{Recovery ladder.} A failed warm solve (non fail-fast) retries as
    a full \code{"seed"} adaptive solve (\code{:906}); the subprocess path adds a
    one-shot reseed (\code{run\_gain\_map.py:424}).
\end{enumerate}

\paragraph{Where the traversal stops: the harmonic-balance fold.}
As pump power rises a column reaches a \emph{turning point} (fold) of
Eq.~\ref{eq:residual}: above it the branch has no nearby solution and Newton
cannot re-converge. The code does not attempt to pass this fold. Instead it
detects repeated failures and culls the rest of the column:
\code{--fold-skip-patience} (default 2) marks all higher-power cells
\code{SKIP\_PAST\_FOLD} (gain \code{NaN}) without solving
(\code{past\_fold\_skip\_row}, \code{run\_gain\_map.py:826}). This is a known
sharp edge: at broadband operating points the true solution can need
$\sim$15 Newton iterations, and a too-tight Newton cap plus
\code{fold-skip-patience 2} can \emph{false-trip} and skip the genuine broadband
region before the real fold.\footnote{Documented internally: the false-trip
culls the 2c broadband operating point; mitigations are looser Newton caps and
higher fold-skip patience.}

\paragraph{One important structural fact.}
Warm-starting is \emph{purely one-dimensional}: it chains \emph{within} a
frequency column along the power axis. Each new frequency column is restarted
from the linear phasor seed --- there is \emph{no} warm start from the adjacent
converged \emph{column}. The 2-D grid is thus traversed as a set of independent
1-D sweeps. This is the single biggest lever the alternatives in
\S\ref{sec:alternatives} act on.

\subsection{Code map}
\begin{center}
\small
\begin{tabular}{@{}p{4.6cm}p{4.2cm}p{5.2cm}@{}}
\toprule
\textbf{Concept} & \textbf{Location} & \textbf{Role} \\
\midrule
Fixed-step continuation & \code{solver.py:402} & cold reference, $N{=}20$ $\lambda$ ladder \\
Adaptive-step continuation & \code{solver.py:463} & step grow/shrink + fixed fallback \\
Single warm Newton solve & \code{solver.py:383} & inter-cell warm start at $\lambda{=}1$ \\
Newton--Krylov corrector & \code{solver.py:101} & GMRES on true Jacobian, precond. \\
Secant predictor (in $\lambda$) & \code{solver.py:422,\,501} & 1st-order extrapolation in continuation \\
Linear phasor seed & \code{seeds.py:46} & homotopy start $\mathcal{L}_1 \Xv_1 = \mathbf{S}_1$ \\
Warm pass (2-D traversal) & \code{run\_gain\_map.py:849} & column-major, low$\to$high power \\
Secant predictor (in $I_p$) & \code{run\_gain\_map.py:785} & inter-cell power-axis extrapolation \\
Fold short-circuit & \code{run\_gain\_map.py:826} & \code{SKIP\_PAST\_FOLD} culling \\
\bottomrule
\end{tabular}
\end{center}

% ===========================================================================
\section{Why the fold is the crux}
% ===========================================================================
Natural-parameter continuation parameterizes the solution branch by the physical
parameter (here pump amplitude). At a fold the branch is \emph{vertical} in that
parameter: $\partial \Xv / \partial \lambda \to \infty$ and the Jacobian
$\Jac$ becomes singular. Any method that insists on advancing the parameter
monotonically \emph{must} stall there --- which is exactly why the current code
treats the fold as a hard wall and culls beyond it. The literature's central
message is that this wall is an artefact of the \emph{parameterization}, not of
the physics: reparameterizing the branch (by arclength) lets a solver round the
fold and continue onto the upper branch \cite{Keller1977,AllgowerGeorg2003}.
Whether the upper (post-fold) branch is physically wanted for a TWPA gain map is
a separate design question; but even if only the lower branch matters, a
fold-aware method locates the fold precisely instead of bracketing it with a
heuristic patience counter.

% ===========================================================================
\section{Alternatives from the literature}
\label{sec:alternatives}
% ===========================================================================

\subsection{Pseudo-arclength continuation (round the fold)}
Instead of fixing $\lambda$ and solving for $\Xv$, treat \emph{both} $\Xv$ and
$\lambda$ as unknowns and add one scalar constraint that advances a
\emph{pseudo-arclength} $s$ along the branch \cite{Keller1977}:
\begin{equation}
  \begin{cases}
    \Res(\Xv,\lambda) = 0,\\[2pt]
    \dot{\Xv}_0^{\!\top}(\Xv - \Xv_0) + \dot{\lambda}_0(\lambda - \lambda_0) - \Delta s = 0,
  \end{cases}
  \label{eq:pal}
\end{equation}
where $(\dot{\Xv}_0,\dot{\lambda}_0)$ is the unit tangent at the current point.
The augmented $(N{+}1)\times(N{+}1)$ system stays non-singular \emph{through} the
fold because the branch is never vertical in $s$. This is the standard way HB
frequency-response and force-response curves are traced past turning points in
mechanics and RF/microwave engineering
\cite{ArclengthHB_Rotor,AllgowerGeorg2003}. Cost: the augmented system and the
tangent solve; benefit: no fold heuristic, and the fold power is found exactly.

\subsection{Tangent (Euler) predictor vs.\ the current secant}
The code's predictors are trivial (copy) and secant
(Eq.~\ref{eq:secant}). The next step up is the \emph{tangent / Euler predictor}:
solve $\Jac\,\dot{\Xv} = \partial\Res/\partial\lambda$ once to get the true
branch tangent and predict $\Xv_{\text{guess}} = \Xv_0 + \Delta\lambda\,
\dot{\Xv}_0$ \cite{AllgowerGeorg2003}. It is first-order-accurate like the
secant but uses derivative information rather than a finite difference of past
points, so it is robust when the previous two points are far apart or the branch
curves sharply --- precisely the near-fold regime. \code{BifurcationKit.jl}'s
predictor page is a compact catalogue of secant/tangent/Bordered/polynomial
predictors and their trade-offs \cite{BifurcationKitPred}.

\subsection{Affine-invariant adaptive step control}
The current adaptive controller uses fixed $\times1.5/\times0.5$ heuristics.
Deuflhard's affine-invariant Newton theory gives a principled step-length and
damping controller driven by the observed Newton contraction, with convergence
monitors that resize the continuation step from the corrector's behaviour rather
than a fixed ratio \cite{Deuflhard2004}. This tends to take large steps on the
flat part of a column and automatically shorten them approaching the fold.

\subsection{Pseudo-transient continuation ($\Psi$tc)}
For a single stiff cell, pseudo-transient continuation embeds the root-find in a
fake time march, $\frac{\Xv^{n+1}-\Xv^{n}}{\delta_n} + \Res(\Xv^{n+1}) = 0$,
growing the pseudo-timestep $\delta_n$ (SER rule) as the residual falls, so early
iterations are globally stable and late iterations recover Newton's quadratic
rate \cite{KelleyKeyes1998,Kelley2003}. This is an alternative to the
$\lambda$-ladder for getting a hard cell to converge from a poor guess, and it
degrades gracefully near the fold (it slows rather than diverges), which could
replace the stall/deadline aborts.

\subsection{2-D map-traversal strategies (the biggest lever)}
Today the map is a bundle of independent 1-D power sweeps
(\S\ref{sec:intercell}); the frequency axis carries \emph{no} warm information.
Alternatives that use the second dimension:
\begin{itemize}[leftmargin=*]
  \item \textbf{Nearest-neighbour seeding on the grid.} Seed each new cell from
    the \emph{best already-converged neighbour} in \emph{either} axis (power or
    frequency), not only the cell below it. A converged cell in the adjacent
    frequency column is often a far better guess than a cold linear-phasor seed
    at the bottom of a new column.
  \item \textbf{Spanning-tree / flood-fill traversal.} Treat converged cells as a
    growing front and always solve the frontier cell whose neighbour guess is
    cheapest (a Dijkstra/Prim-like order over the grid), rather than a fixed
    column-major order. This maximizes the number of cells reached by a single
    warm Newton solve.
  \item \textbf{Boustrophedon (serpentine) ordering.} Traverse columns in
    alternating power direction so the end of one column is adjacent (in power)
    to the start of the next, enabling a warm hand-off \emph{between} columns
    instead of a cold restart.
  \item \textbf{Two-parameter / manifold continuation.} Fold detection can be
    promoted to \emph{fold-following}: continue the fold curve itself in the
    $(f_p, P)$ plane as a codimension-1 manifold, mapping the exact boundary of
    the solvable region in one sweep instead of re-detecting it per column
    \cite{DhoogeMatcont2003,AllgowerGeorg2003}.
\end{itemize}
These change only the \emph{order} of solves and the \emph{source} of each
initial guess; the per-cell physics and corrector are untouched, so they compose
with everything in \S\ref{sec:intracell}.

\subsection{Deflated continuation (find multiple branches)}
If both the lower and upper (post-fold) HB branches --- or genuinely distinct
coexisting steady states --- are of interest, \emph{deflated continuation}
systematically discovers them. After finding one solution, the residual is
multiplied by a deflation operator that makes Newton \emph{repelled} from the
known root, so re-solving from the same guess converges to a \emph{different}
solution; sweeping the parameter then traces disconnected branches without a
priori tangent information \cite{FarrellBirkissonFunke2015}. This is the tool of
choice when a map cell may have more than one physical steady state.

\subsection{Anderson / vector-extrapolation acceleration}
Orthogonal to the above: Anderson acceleration mixes several previous residuals
to accelerate the fixed-point/corrector iteration, often cutting iteration counts
on mildly nonlinear cells for negligible cost and memory
\cite{WalkerNi2011}. It can wrap the existing Newton--Krylov corrector rather
than replace it.

\subsection{Domain context: HB continuation is standard practice}
Arclength/pseudo-arclength continuation wrapped around a harmonic-balance
residual is the established way to compute nonlinear frequency- and
amplitude-response curves with turning points, both in structural dynamics
\cite{ArclengthHB_Rotor} and in nonlinear microwave/RF circuit simulators;
JTWPA-specific modelling increasingly leans on commercial and custom HB solvers
for exactly the pump-depletion / multi-harmonic regime this map probes
\cite{JTWPA_CircuitSim2022,JTWPA_EMcosim2024}. The methods above are therefore
well-trodden rather than exotic.

% ===========================================================================
\section{Summary comparison}
% ===========================================================================
\begin{center}
\small
\begin{tabular}{@{}p{3.5cm}p{2.4cm}p{4.0cm}p{4.3cm}@{}}
\toprule
\textbf{Method} & \textbf{Scope} & \textbf{Passes the fold?} & \textbf{Relation to current code} \\
\midrule
Fixed / adaptive $\lambda$ continuation & intra-cell & No (stalls at fold) & \emph{in use} \\
Secant predictor & intra/inter-cell & No & \emph{in use} \\
Warm-start single Newton & inter-cell & No & \emph{in use} \\
Tangent (Euler) predictor & intra/inter & No, but robust near it & drop-in predictor upgrade \\
Pseudo-arclength & intra-cell & \textbf{Yes} & augments residual with Eq.~\ref{eq:pal} \\
Affine-invariant step control & intra-cell & sharper approach & replaces $\times1.5/\times0.5$ \\
Pseudo-transient ($\Psi$tc) & per-cell solve & graceful slow-down & alt.\ to $\lambda$-ladder \\
Neighbour / tree traversal & inter-cell (2-D) & n/a (ordering) & uses frequency axis for warm starts \\
Fold-following & map-level & maps fold curve & replaces patience heuristic \\
Deflated continuation & per-cell & finds upper branch & discovers extra solutions \\
Anderson acceleration & corrector & n/a & wraps Newton--Krylov \\
\bottomrule
\end{tabular}
\end{center}

% ===========================================================================
\begin{thebibliography}{99}
% ===========================================================================
\bibitem{Keller1977} H.~B.~Keller,
  \emph{Numerical solution of bifurcation and nonlinear eigenvalue problems},
  in \textit{Applications of Bifurcation Theory}, Academic Press, 1977,
  pp.~359--384. (Origin of pseudo-arclength continuation.)

\bibitem{AllgowerGeorg2003} E.~L.~Allgower and K.~Georg,
  \emph{Introduction to Numerical Continuation Methods}, Classics in Applied
  Mathematics 45, SIAM, 2003.
  \url{https://epubs.siam.org/doi/book/10.1137/1.9780898719154}

\bibitem{Deuflhard2004} P.~Deuflhard,
  \emph{Newton Methods for Nonlinear Problems: Affine Invariance and Adaptive
  Algorithms}, Springer Series in Computational Mathematics 35, 2004.

\bibitem{KelleyKeyes1998} C.~T.~Kelley and D.~E.~Keyes,
  \emph{Convergence analysis of pseudo-transient continuation},
  SIAM J. Numer. Anal. \textbf{35}(2):508--523, 1998.

\bibitem{Kelley2003} C.~T.~Kelley,
  \emph{Solving Nonlinear Equations with Newton's Method}, SIAM, 2003.

\bibitem{FarrellBirkissonFunke2015} P.~E.~Farrell, {\'A}.~Birkisson, and
  S.~W.~Funke, \emph{Deflation techniques for finding distinct solutions of
  nonlinear partial differential equations}, SIAM J. Sci. Comput.
  \textbf{37}(4):A2026--A2045, 2015.
  \url{https://epubs.siam.org/doi/10.1137/140984798}

\bibitem{WalkerNi2011} H.~F.~Walker and P.~Ni,
  \emph{Anderson acceleration for fixed-point iterations},
  SIAM J. Numer. Anal. \textbf{49}(4):1715--1735, 2011.

\bibitem{DhoogeMatcont2003} A.~Dhooge, W.~Govaerts, and Yu.~A.~Kuznetsov,
  \emph{MATCONT: A MATLAB package for numerical bifurcation analysis of ODEs},
  ACM Trans. Math. Software \textbf{29}(2):141--164, 2003.

\bibitem{BifurcationKitPred} R.~Veltz et al.,
  \emph{BifurcationKit.jl --- Predictors / correctors}, documentation.
  \url{https://bifurcationkit.github.io/BifurcationKitDocs.jl/dev/Predictors/}

\bibitem{ArclengthHB_Rotor} P.~Sundararajan and S.~T.~Noah,
  \emph{Harmonic balance method with arc-length continuation in rotor/stator
  contact problems}, J. Sound and Vibration / ASME (see review in
  \textit{Experimental continuation in nonlinear dynamics}, arXiv:2408.00138).
  \url{https://arxiv.org/abs/2408.00138}

\bibitem{JTWPA_CircuitSim2022} K.~Peng et al.,
  \emph{Simulating the behaviour of travelling-wave superconducting parametric
  amplifiers using a commercial circuit simulator}, Supercond. Sci. Technol.
  \textbf{35}, 2022.
  \url{https://iopscience.iop.org/article/10.1088/1361-6668/ac850b}

\bibitem{JTWPA_EMcosim2024}
  \emph{Modeling Josephson traveling-wave parametric amplifiers with
  electromagnetic and circuit co-simulation}, arXiv:2509.07807.
  \url{https://arxiv.org/abs/2509.07807}
\end{thebibliography}

\end{document}
